\subsection{证明}

\(\mathfrak{G}\) 中的一个示范性文本包括：

\begin{enumerate}
\def\labelenumi{(\arabic{enumi})}
\item
  \(\mathfrak{G}\) 中关系与项的一个辅助形成性构造 .
\item
  \(\mathfrak{G}\) 中的一个证明，即 \(\mathfrak{G}\) 中的一个关系序列，这些关系都出现在辅助形成性构造中，使得对于序列中的每一个关系 \(\pmb{R}\) ，至少满足以下条件之一：
\end{enumerate}

\((a_{1})\boldsymbol{R}\) 是 \(\mathfrak{c}\) 的一条显式公理 .

\((a_{2})\) R 是由将 C 的一个方案应用于辅助形成性构造中出现的项或关系而得出的结果。

\begin{enumerate}
\def\labelenumi{(\alph{enumi})}
\setcounter{enumi}{1}
\tightlist
\item
  有两个关系 \(S, T\) 在序列中位于 \(R\) 之前，使得 \(T\) 是 \(S \Longrightarrow R\) .
\end{enumerate}

¶ \(\mathfrak{G}\) 中的一个定理，是在 \(\mathfrak{G}\) 的一个证明中出现的关系。

因此，这一概念本质上依赖于所考虑理论在被描述时的状态。理论 \(\mathfrak{G}\) 中的一个关系，当人们成功地将它插入到 \(\mathfrak{G}\) 的一个证明中时，便成为 \(\mathfrak{G}\) 中的一条定理。要说 \(\mathfrak{G}\) 中的一个关系``不是 \(\mathfrak{G}\) 中的定理''，若不参照理论 \(\mathfrak{G}\) 的发展阶段，便不具有任何意义。

\(\mathfrak{G}\) 中的定理也被称为“\(\mathfrak{G}\) 中的真关系”（或“命题”、“引理”、“推论”等）。设 \(R\) 是 \(\mathfrak{G}\) 中的一个关系，设 \(x\) 是一个字母，\(T\) 是 \(\mathfrak{G}\) 中的一个项；若 \((T|x)R\) 是 \(\mathfrak{G}\) 中的定理，则称 \(T\) 满足关系 \(R\) （在 \(\mathfrak{G}\) 中；或称为 \(R\) 的解），这里 \(R\) 被视为 \(x\) 中的关系。

一个关系，若其否定是 \(\mathfrak{c}\) 中的定理，则被称为在 \(\mathfrak{c}\) 中是假的。一个理论 \(\mathfrak{c}\) ，当其中写有一个在 \(\mathfrak{c}\) 中既真又假的关系时，就被称为矛盾的。

这里我们再次涉及一个依赖于理论特定发展状态的概念。读者应当警惕一种混淆（不幸的是，这种混淆由“假”一词的直观含义所暗示），即认为一旦证明了关系 R 在 \(\mathfrak{G}\) 中为假，就由此确立了 R 在 \(\mathfrak{G}\) 中“不真”（严格来说，正如我们上面所指出的，这个短语在数学中没有精确的含义）。

¶ 在下文中，我们将给出称为演绎准则的元数学准则，它们将使我们能够缩短证明。这些准则将用字母 C 后跟一个数字来表示。

C1（三段论）。设 A 和 B 是理论 C 中的关系。如果 A 和 \(A \Longrightarrow B\) 是 C 中的定理，则 B 是 C 中的定理。

设 \(R_1, R_2, \ldots, R_n\) 是 \(\mathfrak{G}\) 中的一个证明，在其中 \(A\) 出现，并且令 \(S_1, S_2, \ldots, S_p\) 是 \(\mathfrak{G}\) 中的一个证明，在其中 \(A \Longrightarrow B\) 出现。显然 \(R_1, R_2, \ldots, R_n, S_1, S_2, \ldots, S_p\) 是 \(\mathfrak{G}\) 中的一个证明，在其中 \(A\) 和 \(A \Longrightarrow B\) 都出现。因此

\[
\boldsymbol {R} _ {1}, \boldsymbol {R} _ {2}, \dots , \boldsymbol {R} _ {n}, \boldsymbol {S} _ {1}, \boldsymbol {S} _ {2}, \dots , \boldsymbol {S} _ {p}, \boldsymbol {B}
\]

是 \(\mathfrak{c}\) 中的一个证明，从而 \(\pmb{B}\) 是 \(\mathfrak{c}\) 中的定理 .

